\documentclass[output=paper]{langscibook}
\ChapterDOI{10.5281/zenodo.22880300}
\title{\textit{Paka suva} – Fish poisoning leaves}
\author{Helen Kobaa Magum,  Shalom Magum, and Joyce Maion}
\abstract{}
\begin{document}
\maketitle\label{chap:36}
\begin{figure}[H]
    \includegraphics[trim=0 0 0 0,clip,width=.6\linewidth]{figures/36_1.jpg}
    \caption{\textit{Paka suva} -- Fish poisoning leaves}
    \label{fig:36_1}
\end{figure}
\begin{xltabular}{\textwidth}{lXX}
	(\ref{ex:text29-1})&\textit{Enaa kahi vahutate ni nom havee to suvasuva vaavaha mi raara a paka suva, eara a ta Teapu.}&I am going to describe how we poison fish with \textit{suva} leaves, we the Teop people.\\
	(\ref{ex:text29-2})&\textit{Beara mei rara tea nao tea suvasuva tea masi,}&Before we go to poison fish on the reef at low tide,\\
	(\ref{ex:text29-3})&\textit{eara toro matatopo roho amaa taba to pasi dee raara tea suvasuva.}&we must prepare the things that we take with us to poison fish.\\
	(\ref{ex:text29-4})&\textit{Momohu eara toro nao rasuu tea osi paka suva teo hum to neba komana raara eve.}&First we must go to the jungle to pick the \textit{suva} leaves at the place where we plant them.\\
	(\ref{ex:text29-5})&\textit{Beara osi, eara toro osi vaberobero, eara repaa dee voosu eve.}&When we pick the leaves, we must pick plenty and carry them home.\\
	(\ref{ex:text29-6})&\textit{Beara tau nao, eara repaa vaobete a paka suva teo hoi.}&Before we go, we put the \textit{suva} leaves into a basket.\\
	(\ref{ex:text29-7})&\textit{Eara repaa kae vo tea masi.}&Then we carry it to the reef.\\
	(\ref{ex:text29-8})&\textit{Beara no suguna tea kopua, eara repaa tara voen, o iana na tei rori teebona,}&When we arrive at a pool and see that there are the fish,\\
	(\ref{ex:text29-9})&\textit{eara repaa vahio o hoi paka suva.}&we put the basket with the \textit{suva} leaves down.\\
	(\ref{ex:text29-10})&\textit{Eara repaa taneo tea buhaha a paka suva.}&Then we start rubbing the \textit{suva} leaves into pulp with our hands.\\
	(\ref{ex:text29-11})&\textit{Buhaha bata mi raara bono kasuana to karusu bata raara komana kopua.}&We rub them with sand that we collect with our hands in the pool.\\
	(\ref{ex:text29-12})&\textit{Eara buhaha batara, beara peo binubinu bata ni rae komana kopua, ge tea maa avaava teo vasu.}&We rub them into pulp, while scattering them everywhere in the pool or in the caves of the stones.\\
	(\ref{ex:text29-13})&\textit{Beara suva batara, eara repaa tara voen o iana na mate batari, eara repaa antee tea buhaha a paka suva.}&When we do the poisoning, and we see that the fish are dying, we stop rubbing the \textit{suva} leaves into pulp.\\
	(\ref{ex:text29-14})&\textit{Eara repaa taneo tea mosi o iana to bubuutoo bata maana.}&Then we start killing the struggling fish with a knife.\\
	(\ref{ex:text29-15})&\textit{A paka suva vai to paa buhaha vuru ara, na kaokahi koa raara komana kopua,}&The \textit{suva} leaves that we have rubbed into pulp we leave in the pool,\\
	(\ref{ex:text29-16})&\textit{ei kou paa kavara vaevuru a kikis nae.}&because their strength has already gone.\\
	(\ref{ex:text29-17})&\textit{Eara repaa kae voosu o hoi iana.}&Then we carry the basket full of fish home.\\
	(\ref{ex:text29-18})&\textit{Eve koai o sii vahuatate kakato vaa tea suvasuva nia paka suva.}&This is the short description of poisoning fish with \textit{suva} leaves.\\
\end{xltabular}
\end{document}